\lysh{14375}{4}

\begin{alist}
\item Η διακρίνουσα του τριωνύμου είναι:
$\Delta = (-\mu)^2 - 4 \cdot 1 \cdot (-2) = \mu^2 + 8 > 0$ για κάθε $\mu \in \mathbb{R}$.
Άρα η εξίσωση $f(x) = 0$ έχει δύο ρίζες πραγματικές άνισες για κάθε $\mu \in \mathbb{R}$.

\item Είναι:
\begin{gather*}
f(-2) = 4 + 2\mu - 2 = 2\mu + 2, \\
f(1) = 1 - \mu - 2 = -\mu - 1, \\
f(3) = 9 - 3\mu - 2 = 7 - 3\mu.
\end{gather*}
Για να βρίσκονται τα $x = -2$ και $x = 3$ εκτός του διαστήματος των ριζών της εξίσωσης
$f(x) = 0$ ενώ το $x = 1$ εντός του διαστήματος των ριζών της εξίσωσης $f(x) = 0$, θα 
πρέπει να ισχύει:
\[\begin{cases} f(-2) > 0 \\ f(1) < 0 \\ f(3) > 0 \end{cases} \Leftrightarrow \begin{cases} 2\mu + 2 > 0 \\ -\mu - 1 < 0 \\ 7 - 3\mu > 0 \end{cases} \Leftrightarrow \begin{cases} \mu > -1 \\ \mu > -1 \\ \mu < \dfrac{7}{3} \end{cases} \Leftrightarrow -1 < \mu < \dfrac{7}{3}.\]

\item 
\begin{rlist}
\item Για να είναι τα $f(-2), f(1), f(3)$ με τη σειρά που δίνονται διαδοχικοί όροι
γεωμετρικής προόδου θα πρέπει να ισχύει:
\begin{align*}
(f(1))^2 = f(-2) \cdot f(3) \ &\Leftrightarrow \\
(-\mu - 1)^2 = (2\mu + 2) \cdot (7 - 3\mu) \ &\Leftrightarrow \\
(\mu + 1)^2 - 2(\mu + 1) \cdot (7 - 3\mu) = 0 \ &\Leftrightarrow \\
(\mu + 1) \cdot (\mu + 1 - 14 + 6\mu) = 0 \ &\Leftrightarrow \\
(\mu + 1) \cdot (7\mu - 13) = 0 \ &\Leftrightarrow \\
\mu = -1 \text{ ή } 7\mu - 13 = 0 \ &\Leftrightarrow \\
\mu &= -1 \text{ ή } \mu = \dfrac{13}{7}.
\end{align*}
Άρα $\mu = \dfrac{13}{7}$ αφού $\mu \in \left(-1, \dfrac{7}{3}\right)$.

\item Για $\mu = \dfrac{13}{7}$ είναι:
\begin{gather*}
f(-2) = 2 \cdot \dfrac{13}{7} + 2 = \dfrac{40}{7}, \\
f(1) = -\dfrac{13}{7} - 1 = -\dfrac{20}{7}\text{ και} \\
f(3) = 7 - 3 \cdot \dfrac{13}{7} = \dfrac{10}{7}.
\end{gather*}
Άρα ο λόγος $\lambda = \dfrac{-\dfrac{20}{7}}{\dfrac{40}{7}} = -\dfrac{1}{2}$.
\end{rlist}
\end{alist}

